% ---
% id: statistics/thm:box-oracle
% kind: thm
% title: "Acceptance support and the exact box oracle"
% summary: >-
%   Over a non-empty canonical box, filling the coordinates of smallest
%   coefficient first minimises a rational linear functional exactly, with a
%   dual witness of the same value. The outcomes with a positive upper bound
%   include every outcome on which a point of the box puts weight, and can
%   include more.
% prominence: main
% status: implemented
% uses: [def:taxonomy, thm:geometry-validation, def:matched-alphabet]
% mentions: [meat-core/frank-wolfe/def:first-order-setting, meat-core/frank-wolfe/rem:acc-side, foundations/req:rounding-direction]
% assumes: [A-ALPHABET]
% refs: []
% provenance: newly-derived
% implements: certiqs_key.statistics.minimization:box_minimum
% lean:
% ---

The consuming engine minimises linear functionals over the acceptance set and
needs to know on which outcomes an accepted frequency can put weight. Over a
canonical box the minimum has a closed form, and the outcomes with a positive
upper bound form a closed-form superset of the outcomes that can carry weight.
The dual witness of (iii) has the format of instantiation~(b) of
\xref{meat-core/frank-wolfe/rem:acc-side}: \eqref{eq:dualfeas} implies, at
every outcome, the dual-feasibility inequality that instantiation requires,
and the witness
attains the minimum, \eqref{eq:dualvalue}, which is why the bound it gives is
tight. Rounded down, that bound remains a lower bound, the direction of
\xref{foundations/req:rounding-direction}, so the remark's certified
acceptance-side bound $\bacc$ may be taken as any downward-rounded evaluation
of $\langle g,\mathbf{q}^{\mathrm{opt}}\rangle$.

\begin{theorem}\label{thm:box-oracle}
Let $\Chatset$ and $n$ be a \tref{matched alphabet and round count}
(Definition~\ref{def:matched-alphabet}), let $\acSet$ be a non-empty
\tref[canonical box form]{canonical box}, \eqref{eq:boxset}, with
$0\le\qlo\le\qhi$, and let $g\in\mathbb{Q}^{\Chatset}$.
\begin{enumerate}[label=(\roman*),leftmargin=2.2em]
\item \emph{(acceptance support)} The set
$\Chatset^{+}\coloneqq\{\hat c\in\Chatset: q^{+}(\hat c)>0\}$ contains
$\operatorname{supp}\mathbf{q}$ for every $\mathbf{q}\in\acSet$; hence it
contains the acceptance support
$\Chacc=\bigcup_{\mathbf{q}\in\acSet}\operatorname{supp}\mathbf{q}$ of
\xref{meat-core/frank-wolfe/def:first-order-setting}, which it need not equal: the box
constraints together with the simplex constraint can force $q(\hat c)=0$ at an
outcome with $q^{+}(\hat c)>0$. For a $\mathrm{singleton}$ geometry, which
declares no box, the union of the supports is
$\operatorname{supp}\mathbf{q}^{\ast}$, the acceptance support itself.
\item \emph{(exact minimum in closed form)} Enumerate
$\Chatset=\{c_1,\dots,c_m\}$ so that $g(c_1)\le\dots\le g(c_m)$. Put
$r\coloneqq1-\sum_{\hat c}q^{-}(\hat c)$ and
$D_i\coloneqq q^{+}(c_i)-q^{-}(c_i)\ge0$, let $j^{\ast}$ be the least index
with $\sum_{i\le j^{\ast}}D_i\ge r$, and define
\begin{equation}\label{eq:greedy}
  q^{\mathrm{opt}}(c_i)\coloneqq
  \begin{cases}
    q^{+}(c_i), & i<j^{\ast},\\[2pt]
    q^{-}(c_i)+\big(r-\sum_{i'<j^{\ast}}D_{i'}\big), & i=j^{\ast},\\[2pt]
    q^{-}(c_i), & i>j^{\ast}.
  \end{cases}
\end{equation}
Then $j^{\ast}$ exists, $\mathbf{q}^{\mathrm{opt}}\in\acSet$, and
\begin{equation}\label{eq:exactmin}
  \langle g,\mathbf{q}^{\mathrm{opt}}\rangle
  \;=\;\min_{\mathbf{q}\in\acSet}\ \langle g,\mathbf{q}\rangle ,
\end{equation}
the minimum being attained. All quantities are exact rationals when
$g,\qlo,\qhi$ are.
\item \emph{(dual witness)} With $B,\mathbf{b}$ as in
\eqref{eq:boxrows}, $\lambda\coloneqq g(c_{j^{\ast}})$, $v\coloneqq-\lambda$ and
$\mathbf{u}\coloneqq(\mathbf{u}^{+},\mathbf{u}^{-})$ given by
$u^{+}(\hat c)\coloneqq\max\{\lambda-g(\hat c),0\}$ and
$u^{-}(\hat c)\coloneqq\max\{g(\hat c)-\lambda,0\}$, one has
$\mathbf{u}\ge0$,
\begin{equation}\label{eq:dualfeas}
  g(\hat c)+(B^{\transpose}\mathbf{u})(\hat c)+v=0
  \qquad\text{for every }\hat c\in\Chatset,
\end{equation}
in particular the left-hand side is $\ge0$, and
\begin{equation}\label{eq:dualvalue}
  -\langle\mathbf{b},\mathbf{u}\rangle-v
  \;=\;\langle g,\mathbf{q}^{\mathrm{opt}}\rangle
  \;=\;\min_{\mathbf{q}\in\acSet}\langle g,\mathbf{q}\rangle .
\end{equation}
The lower bound $-\langle\mathbf{b},\mathbf{u}\rangle-v$ on the minimum that
this witness gives is therefore \emph{tight} for the canonical box, and every
downward-rounded evaluation of $\langle g,\mathbf{q}^{\mathrm{opt}}\rangle$ is
a lower bound on $\min_{\mathbf{q}\in\acSet}\langle g,\mathbf{q}\rangle$.
\end{enumerate}
\end{theorem}

\begin{proof}
\emph{(i).} If $\mathbf{q}\in\acSet$ and $q^{+}(\hat c)=0$ then
$0\le q(\hat c)\le q^{+}(\hat c)=0$, so $\hat c\notin\operatorname{supp}
\mathbf{q}$; contrapositively $\operatorname{supp}\mathbf{q}\subseteq\Chatset^{+}$,
and taking the union over $\mathbf{q}\in\acSet$ gives
$\Chacc\subseteq\Chatset^{+}$. That the inclusion can be strict: on
$\Chatset=\{c_1,c_2,c_3\}$ with $\qlo=(\tfrac12,\tfrac12,0)$ and
$\qhi=(\tfrac12,\tfrac12,\tfrac1{10})$ the box is non-empty by
\eqref{eq:boxfeas}, the simplex constraint forces $q(c_3)=0$ for every
$\mathbf{q}\in\acSet$, so $\Chacc=\{c_1,c_2\}$ while
$\Chatset^{+}=\{c_1,c_2,c_3\}$. For a singleton the union
$\bigcup_{\mathbf{q}\in\acSet}\operatorname{supp}\mathbf{q}$ is
$\operatorname{supp}\mathbf{q}^{\ast}$ by inspection, which is therefore the
acceptance support $\Chacc$.

\emph{(ii).} \emph{Existence of $j^{\ast}$.} $\acSet\neq\varnothing$ gives
$\sum_{\hat c}q^{+}(\hat c)\ge1$ by Theorem~\ref{thm:geometry-validation}\,(iii), hence
$\sum_{i\le m}D_i=\sum_{\hat c}q^{+}(\hat c)-\sum_{\hat c}q^{-}(\hat c)\ge
1-\sum_{\hat c}q^{-}(\hat c)=r$, so the index set is non-empty and has a least
element. Also $r\ge0$, again by Theorem~\ref{thm:geometry-validation}\,(iii).

\emph{Feasibility of $\mathbf{q}^{\mathrm{opt}}$.} For $i\neq j^{\ast}$ the value is
one of the two box endpoints. At $i=j^{\ast}$ the increment
$w\coloneqq r-\sum_{i'<j^{\ast}}D_{i'}$ satisfies $w>0$ when
$j^{\ast}>1$, by minimality of $j^{\ast}$, and $w=r\ge0$ when
$j^{\ast}=1$; and $w\le D_{j^{\ast}}$ by the defining inequality
$\sum_{i\le j^{\ast}}D_i\ge r$. Hence
$q^{-}(c_{j^{\ast}})\le q^{\mathrm{opt}}(c_{j^{\ast}})\le q^{+}(c_{j^{\ast}})$, and
$\mathbf{q}^{\mathrm{opt}}\ge\qlo\ge0$. Summing,
$\sum_i q^{\mathrm{opt}}(c_i)=\sum_i q^{-}(c_i)+\sum_{i<j^{\ast}}D_i+w
=\sum_i q^{-}(c_i)+r=1$. So $\mathbf{q}^{\mathrm{opt}}\in\acSet$.

\emph{Optimality.} Let $\mathbf{q}\in\acSet$ and $\lambda=g(c_{j^{\ast}})$.
Since $\sum_i\big(q(c_i)-q^{\mathrm{opt}}(c_i)\big)=1-1=0$,
\begin{equation*}
  \langle g,\mathbf{q}\rangle-\langle g,\mathbf{q}^{\mathrm{opt}}\rangle
  =\sum_{i=1}^{m}\big(g(c_i)-\lambda\big)
   \big(q(c_i)-q^{\mathrm{opt}}(c_i)\big).
\end{equation*}
Each summand is non-negative: for $i<j^{\ast}$ the ordering gives
$g(c_i)-\lambda\le0$ while $q(c_i)-q^{\mathrm{opt}}(c_i)=q(c_i)-q^{+}(c_i)\le0$; for
$i>j^{\ast}$, $g(c_i)-\lambda\ge0$ while
$q(c_i)-q^{\mathrm{opt}}(c_i)=q(c_i)-q^{-}(c_i)\ge0$; and at $i=j^{\ast}$ the first
factor is $0$. Hence
$\langle g,\mathbf{q}\rangle\ge\langle g,\mathbf{q}^{\mathrm{opt}}\rangle$ for every
feasible $\mathbf{q}$, which is \eqref{eq:exactmin} with attainment at
$\mathbf{q}^{\mathrm{opt}}$.

\emph{(iii).} Non-negativity of $\mathbf{u}$ is clear. With the row structure
\eqref{eq:boxrows}, $B^{\transpose}\mathbf{u}=\mathbf{u}^{+}-\mathbf{u}^{-}$,
and for every real $x$ one has $\max\{x,0\}-\max\{-x,0\}=x$; applying this at
$x=\lambda-g(\hat c)$ gives
$u^{+}(\hat c)-u^{-}(\hat c)=\lambda-g(\hat c)$, so
$g(\hat c)+(B^{\transpose}\mathbf{u})(\hat c)+v
=g(\hat c)+\lambda-g(\hat c)-\lambda=0$, which is \eqref{eq:dualfeas}.

For \eqref{eq:dualvalue}, note first that the sorted order and
$\lambda=g(c_{j^{\ast}})$ give $u^{-}(c_i)=0$ for $i\le j^{\ast}$ and
$u^{+}(c_i)=0$ for $i\ge j^{\ast}$. Then
\begin{align*}
  -\langle\mathbf{b},\mathbf{u}\rangle-v
  &=-\langle\qhi,\mathbf{u}^{+}\rangle+\langle\qlo,\mathbf{u}^{-}\rangle
    +\lambda\\
  &=-\sum_{i<j^{\ast}}q^{+}(c_i)\big(\lambda-g(c_i)\big)
    +\sum_{i>j^{\ast}}q^{-}(c_i)\big(g(c_i)-\lambda\big)+\lambda\\
  &=\lambda\Big(1-\sum_{i<j^{\ast}}q^{+}(c_i)-\sum_{i>j^{\ast}}q^{-}(c_i)\Big)
    +\sum_{i<j^{\ast}}g(c_i)q^{+}(c_i)+\sum_{i>j^{\ast}}g(c_i)q^{-}(c_i)\\
  &=\lambda\,q^{\mathrm{opt}}(c_{j^{\ast}})+\sum_{i\neq j^{\ast}}g(c_i)q^{\mathrm{opt}}(c_i)
   \;=\;\langle g,\mathbf{q}^{\mathrm{opt}}\rangle ,
\end{align*}
where the fourth line uses $q^{\mathrm{opt}}(c_i)=q^{+}(c_i)$ for $i<j^{\ast}$,
$q^{\mathrm{opt}}(c_i)=q^{-}(c_i)$ for $i>j^{\ast}$, $\sum_i q^{\mathrm{opt}}(c_i)=1$ and
$\lambda=g(c_{j^{\ast}})$. Combining with \eqref{eq:exactmin} gives
\eqref{eq:dualvalue}. In particular $-\langle\mathbf{b},\mathbf{u}\rangle-v$
is a lower bound on the minimum that it equals, and a downward-rounded
evaluation of $\langle g,\mathbf{q}^{\mathrm{opt}}\rangle$ is at most
$\langle g,\mathbf{q}^{\mathrm{opt}}\rangle$, which is the minimum by
\eqref{eq:exactmin}.
\end{proof}
